% Body of the thesis summary in English: content only, without a preamble or title page.
% Included by summary_en.tex, which provides the document class, style, and title page.
% Kept separate so that it can also be included in ../thesis/thesis.tex
% if the summary is to become part of the thesis.

\section{Introduction}

The growing use of sensors, cameras, and connected devices makes sending all raw data to a remote server increasingly inefficient and, in some cases, infeasible. Applications such as visual inspection on a production line, robotic arm guidance, and environmental monitoring often require a response with very little delay, or low latency, even when network connectivity is unavailable or unreliable. \emph{Edge computing} addresses this need by performing at least part of the processing on the device that generates the data or in its immediate vicinity.

Edge computing has already moved beyond the experimental stage. According to a forecast published by IDC in 2026, worldwide spending on edge computing is expected to rise from \$265 billion in 2025 to \$450 billion in 2029.\footnote{International Data Corporation, \emph{Edge Computing Global Spending to Grow at 15\%, reaching \$450 Billion by 2029}, 26 February 2026, \url{https://www.idc.com/resource-center/press-releases/edgecomputingforecast/}.} Its strategic importance is also recognised by the European Union, whose Digital Decade programme sets a target of 10,000 climate-neutral, highly secure edge nodes by 2030.\footnote{European Commission, \emph{Cloud computing}, \url{https://digital-strategy.ec.europa.eu/en/policies/cloud-computing}.}

An \emph{edge device} therefore does more than collect and transmit information: it performs enough local processing to produce a usable result. A camera in a car park, for example, can identify vacant spaces and report their number instead of continuously transmitting the entire video stream. The cloud still has a role, providing centralised resources for tasks such as data storage, large-scale analysis, and coordination across multiple devices.

\section{Architectures for Edge Devices}

Edge devices can be built using a range of architectures. One option is to use a microcontroller (MCU) or microprocessor (MPU) within an embedded system, a computing system integrated into a device to perform specific functions. The processor executes software and provides access to hardware peripherals defined by the chip manufacturer. When an application requires an interface or logic that is not available among these standard peripherals, or highly specialised processing, the processor can be paired with an FPGA (\emph{Field-Programmable Gate Array}). Unlike an architecture based solely on an MCU or MPU, an FPGA allows the designer to implement custom hardware logic: a circuit tailored to the task at hand. Regular, repetitive operations can then be accelerated directly in hardware, reducing processing latency and relieving the processor of unnecessary computations and data transfers. This division of work helps improve the energy efficiency of the system, while leaving more processor capacity available for tasks better suited to software, such as decision-making, communication, and logic that changes over time.

The system presented in this thesis was developed on the Terasic DE1-SoC board. Its main device, an Intel Cyclone V SoC, embodies a modern approach to edge device design by combining an FPGA and a processing system based on a dual-core ARM Cortex-A9 MPU on the same chip. Intel refers to the non-reconfigurable part of the chip as the \emph{Hard Processor System} (HPS). It includes the MPU, the DDR3 memory controller, and other peripherals, such as USB controllers and the network interface. Dedicated interconnects allow the MPU to communicate with the logic implemented in the FPGA, which software can access as a peripheral.

\section{Thesis Objectives}

The aim of this thesis is to develop a complete video acquisition and processing system. FPGA logic acquires images from the camera provided for the project, which uses a non-standard interface, while software running on the integrated ARM MPU performs the image processing. The starting point is therefore an unconfigured programmable logic fabric, with neither a camera interface nor a data path for transferring images to system memory.

\section{Video Acquisition Pipeline}

\subsection{Camera Interface}

Developing the system also required work on the physical connection between the camera and the board. At the sensor's data transmission frequency, the jumper wires initially used degraded signal quality to the point where the FPGA could no longer read the data reliably. Two adapter printed circuit boards equipped with active line buffers were therefore designed and built to establish a reliable connection between the camera and the development board.

Once the electrical connection had been established, the sensor's digital interface had to be implemented in the FPGA. The camera, based on the OmniVision OV7670 sensor, outputs pixels over an eight-bit parallel bus, driven by its own timing signal, called the pixel clock, and accompanied by line and frame synchronisation signals. It provides no mechanism for the FPGA to request a pause or a lower transmission rate. The acquisition logic must therefore accept each sample as it arrives, identify the start and end of each frame, and pass the frame data to the transfer pipeline in the correct order.

\subsection{Transferring Images to Memory}

The next step is to move the acquired data into DDR3 system memory for processing. A central requirement is to perform this transfer without using the MPU to copy individual data items, leaving its computational resources available for image processing.

This is achieved using a DMA (\emph{Direct Memory Access}) controller, a hardware component that transfers blocks of data autonomously. The MPU only needs to specify the destination memory address and the transfer size. The DMA controller then handles the incoming stream and writes it to memory without requiring the MPU to execute copy instructions for each byte received. Once the transfer is complete, the controller notifies the MPU through an interrupt, a signal that prompts the processor to execute the corresponding handler.

The processor subsystem of the Cyclone V SoC does not provide a DMA controller ready to perform this particular transfer. Intel does, however, supply a configurable hardware component called \emph{Modular Scatter-Gather DMA} (mSGDMA). In this work, it was integrated into the FPGA logic and connected to the HPS DDR3 memory controller through dedicated interconnects.

\subsection{Synchronisation and Incomplete Frames}

Connecting the acquisition interface to the DMA controller also requires their timing differences to be accommodated. Data arrives under the camera clock, whereas the DMA controller runs under the system clock and may pause briefly because of competition for access to the shared DDR3 memory. Temporary storage elements, or buffers, were therefore inserted between them. These buffers allow data to cross between parts of the circuit driven by different clocks, known as clock domains, and hold incoming data until the transfer can resume.

When downstream stages cannot accept further data, they assert a \emph{backpressure} signal. If the buffers can no longer absorb the incoming data during a stall, this signal propagates back to the acquisition interface. Since the camera cannot be stopped, the logic abandons the current frame and, once the transfer resumes, signals its premature end to the DMA controller. The software also receives the actual number of bytes transferred, allowing it to identify and discard the incomplete image rather than treat it as a valid frame.

\section{Implemented Systems and Demonstration Applications}

The acquisition pipeline described above forms the common foundation of two systems developed in this thesis, named \emph{FPGAlix} and \emph{FPGAsteel}. Both use the FPGA to acquire images and transfer them to memory, while retaining image processing on the MPU, but they adopt different software environments. FPGAlix runs Linux, which provides services and tools for application development. FPGAsteel instead executes the application directly on the processor without an operating system, following a \emph{bare-metal} approach.

\subsection{FPGAlix: Video Acquisition under Linux}

FPGAlix includes a custom kernel driver, a software component that allows Linux to control the acquisition pipeline implemented in the FPGA and expose it to applications through Video4Linux2 (V4L2), the standard Linux interface for video devices. The camera thus appears to applications as an ordinary video device: a program can capture frames without knowing how the hardware and software implement the path to memory. Building on this interface and using open-source libraries available for Linux, a demonstration application called \emph{vision\_core} was developed. It captures the video stream, applies a software filter pipeline running on the MPU, and transmits the result through the board's Ethernet interface. The filters can be reconfigured at runtime without interrupting video processing. The demonstration therefore exercises the entire data path, from the sensor to the FPGA, from memory to the processor, and finally to the network.

\subsection{FPGAsteel: Local Processing and Display}

In FPGAsteel, the absence of an operating system means that peripheral control must be implemented explicitly. Transmitting video as in FPGAlix would have required the development of an Ethernet driver and the necessary communication protocols, a substantial and complex undertaking outside the scope of this thesis. To display the processed video, the system instead uses the DE1-SoC's VGA output, which can be connected directly to a monitor. A second DMA path was therefore implemented in the FPGA alongside the acquisition pipeline. It operates in the opposite direction, reading processed frames from DDR3 memory and sending them to the VGA controller.

Compared with FPGAlix, both DMA controllers in FPGAsteel, one for the camera and one for VGA output, include an additional optional internal module called the \emph{prefetcher}. This module retrieves the buffer address and the other information required for a data transfer directly from a dedicated memory. At the end of each transfer, the prefetcher notifies the MPU through an interrupt; the MPU then selects the next buffer and rearms the DMA controller. The processor thus manages frames without having to transfer individual pixels. The result is an entirely local video path in which the FPGA handles data input and output, leaving the MPU available for image processing.

FPGAsteel was developed alongside a progressive series of demonstration applications. Each introduces and verifies a new function while reusing those already tested: system startup and peripheral control, camera configuration, image transfer to the VGA output, software execution optimisation, startup of the second ARM core, and finally the distribution of processing across both cores. This approach allows individual components to be verified and problems to be identified before integration into the complete system. It also helps less experienced readers become familiar with the work in manageable steps.

The series culminates in the most complex application, which uses both ARM cores to detect a bright object against a dark background. The program converts each frame to greyscale, applies a threshold to separate the object from the background, and uses morphological operations to remove isolated noise and extract the object's contour. It then overlays the contour, a bounding box, and a cross marking the object's centroid on the video displayed on the VGA monitor. This application exercises the complete local acquisition, processing, and display pipeline while placing a measurable load on the MPU.

\subsection{MPU Optimisation and Performance}

Three acceleration features of the ARM cores were used to assess MPU performance. The first is the data cache, a very fast memory close to the processor that exploits locality by storing recently accessed data and data from nearby memory addresses. The second is the instruction cache, which serves a similar purpose by storing program instructions. The third is the set of vector instructions known as NEON, which allow the same operation to be applied to several pixels simultaneously.

The data cache remains enabled in all tests. The final application was compiled in three variants implementing the same algorithm. The baseline uses no NEON vector instructions and runs with the instruction cache disabled; the second variant enables the instruction cache; in the third, the code was explicitly rewritten to use NEON instructions. The corresponding processing times per frame fall from 142 to 47 and then to 17 milliseconds. The fully optimised version is therefore approximately 8.35 times as fast as the baseline.

\section{Conclusions and Future Work}

Starting with an unconfigured FPGA, this work established the complete data path needed to acquire a video stream from a camera with a non-standard interface and transfer its frames to DDR3 memory. The FPGA logic and DMA controller move individual bytes, while the MPU manages transfers at frame level and processes images once they are available in system memory. This infrastructure supports two systems, FPGAlix and FPGAsteel, which provide the same acquisition architecture in Linux and \emph{bare-metal} environments, respectively. The demonstration applications developed for both systems verified the integrated operation of the sensor, FPGA logic, memory, processor, and output interfaces.

Image processing was deliberately kept on the MPU in the demonstration applications to place a substantial load on the processor and DDR3 memory access. This made it possible to observe the behaviour of the acquisition pipeline while the software was analysing frames. The measurements quantify the benefits of the optimisations introduced and provide a reference for future work in which functions may be partitioned between the MPU and FPGA logic.

The resulting platform can serve as a basis for designing an edge device for vision and control. Images can be acquired and processed locally, allowing either processed video to be transmitted or, with an application developed for a specific task, only the extracted information, such as a measurement, a count, or an alert. Processing results can also directly control one or more actuators. On a production line, for example, the device could detect a defective component through visual inspection and activate a mechanism to reject it. Processing close to the sensor reduces both the workload assigned to central systems or the cloud and the volume of data sent over the network. The device can also continue to make decisions when connectivity is unavailable or unreliable. It also avoids delays associated with communication with a remote system. FPGAlix demonstrates the platform in an environment connected through Ethernet, whereas FPGAsteel demonstrates that, in the implemented configuration, the entire acquisition, processing, and display pipeline operates locally, using VGA output and avoiding the delays and timing variability introduced by operating system scheduling, which switches execution between programs.

A further development would be to move regular operations that apply the same computation to each pixel, such as format conversion, thresholding, convolution filtering, and morphological processing, into the FPGA and execute them directly on the incoming video stream. This would reduce MPU load and processing latency, leaving software responsible for higher-level decisions and functions that require more frequent modification. The FPGAlix and FPGAsteel configurations use approximately 15\% and 18\% of the available programmable logic, respectively, and neither uses any of the Cyclone V SoC's 87 hardware blocks dedicated to arithmetic and signal processing. Resources therefore remain available for implementing these accelerators.

The documentation and the hardware and software sources made available in the dedicated repositories, FPGAlix\footnote{\url{https://github.com/tceccarini/FPGAlix}.} and FPGAsteel\footnote{\url{https://github.com/tceccarini/FPGAsteel}.}, describe in detail the steps required to reproduce both systems. This allows future contributors to understand the design choices, reproduce the systems' operation, and use the platform as a foundation for custom applications and further extensions, without having to rebuild the entire acquisition and processing pipeline from scratch.
